\section{Evaluation}

\subsection{Metrics}
\label{sec:metrics}

The ability to assess the quality of the generated videos is a challenging task. While humans can evaluate videos across various visual dimensions \cite{huang2023vbench,bugliarello2024storybench}, we focus primarily on the models' adherence to the provided text and the incorporation of physical commonsense. These are key objectives that conditional generative models must maximize. \hb{We note that several video characteristics such as object motion, video quality, text adherence, physical commonsense, temporal consistency of subject and object etc. are usually intertwined with each other. It is non-trivial to disentangle their effect when humans make decisions. However, focusing on each aspect at a time provides a comprehensive picture of the model capabilities along a specific dimension. In this work, we focus on physical commonsense and semantic adherence.}
\hb{Further, there are diverse ways to acquire human judgments such as dense and sparse feedback. While a dense feedback provides detailed information about the model mistakes, it is hard to acquire and miscalibrated \cite{lozano2008effect,lewis2017user}. Due to the simplicity of binary judgment and its widespread use in text-to-image generative models \cite{li2024aligning,lee2023aligning}, we employ binary feedback (0/1) to evaluate the generated videos in this work. Our experiments will demonstrate that binary feedback effectively highlights differences in the model's quality across various object interactions and levels of task complexity.}

\paragraph{Semantic Adherence (SA).} This metric assesses whether the text caption is semantically grounded in the frames of the generated videos, measuring video-text alignment. Specifically, it assesses if the actions, events, entities, and their relationships are perceived to be correctly depicted in the video frames (e.g., water is flowing into the glass in the generated video for the caption `water pouring into the glass'). In this work, we annotate the generated videos for semantic adherence, denoted as $\text{SA}=\{0, 1\}$. Here, $\text{SA}=1$ indicates that the caption is grounded in the generated video.

\textbf{Physical Commonsense (PC).} This metric evaluates whether the depicted actions, and object's state follow the physics laws in the real-world. For instance, the level of water should increase in the glass as water flows into it, following conversation of mass. In this work, we annotate the physical commonsense of the generated videos, denoted as $\text{PC}=\{0,1\}$. Here, $\text{PC}=1$ indicates that the generated movements and interactions align with intuitive physics that humans acquire with their experience in the real-world. As physical commonsense is entirely grounded in the video, it is independent of the semantic adherence capability of the generated video. In this work, we compute the fraction of the videos for which semantic adherence is high ($\text{SA} = 1$), physical commonsense is high ($\text{PC} = 1$), and joint performance of these metrics is high ($\text{SA}=1, \text{PC} = 1$). 

\subsection{Human Evaluation}
\label{sec:human_evaluation}

We conducted a human evaluation to assess the performance of the generated videos in terms of semantic adherence and physical commonsense using our dataset. The annotations were obtained from a group of qualified Amazon Mechanical Turk (AMT) workers \hb{who were provided with the detailed task description (and clarifications) on a shared slack channel. Subsequently, 14 workers who have studied high-school physics were chosen to perform the annotations after passing a qualification test.} In this task, annotators were presented with a caption and the corresponding generated video without any information about the generative model. They were asked to provide a semantic adherence score (0 or 1) and a physical commonsense score (0 or 1) for each instance. Annotators were instructed to treat semantic adherence and physical commonsense as independent metrics and were shown several solved examples by the authors before starting the main annotation task. In some cases, we find that generative models create static scenes instead of video frames with high motion. Here, we ask annotators to judge the physical plausibility of the static scene in the real world (e.g., a static scene of a folded brick does not follow physical commonsense). If the static scenes are noisy (e.g., unwanted grainy or speckled patterns), we instruct them to consider it as poor physical commonsense. \footnote{The workers were compensated at a rate of $\$18$ per hour.}

The human annotators were not asked to list the violation of the physics laws since it would make the annotations more time-consuming and expensive. Additionally, the current annotations can be performed by annotators experience in the physical world (e.g., workers know that water flows \textit{down} from a tap, shape of a wood log \textit{will not change} while floating on water) instead of advanced education in physics. A screenshot of the human annotation interface is presented in Appendix \ref{sec:human_annotation_screenshot}. 



\subsection{Automatic Evaluation}
\label{sec:automatic_evaluation}

While the human evaluation is more accurate for benchmarking, it is time-consuming and expensive to acquire at scale. \hb{In addition, we want the model developers with limited resources for human evaluation to use our benchmark.} To this end, we design \textbf{\model{}}, a reliable auto-rater for our evaluation dataset. 
% To this end, we evaluate the performance of various zero-shot methods in judging the quality of the generated videos in terms of semantic adherence and physical commonsense.
% Since the previous two models are closed, it is difficult to fine-tune them with custom data. 
Specifically, we use \videocon{}, an open video-text language model with $7$B parameters, that is trained on real videos for robust semantic adherence evaluation \cite{bansal2023videocon}. Specifically, we prompt \videocon{} to generate a text response (\textit{Yes}/\textit{No}) conditioned on the multimodal template. \hb{We provide details about the templates and score computation using \videocon{} in Appendix \ref{app:videocon_details}.} 

Since \videocon{} is not trained on the generated video distribution or equipped to judge physical commonsense, it is not expected to perform well in our setup in a zero-shot manner. To this end, we propose \model{}, an open-source generative video-text model, that can assess the semantic adherence and physical commonsense of the generated videos. Specifically, we finetune \videocon{} by combining the human annotations acquired for the semantic adherence and physical commonsense tasks over the generated videos. \hb{We present the GPT-4Vision \cite{gpt4v} and Gemini-1.5-Pro-Vision \cite{reid2024gemini} baselines in Appendix \ref{app:automatic_eval_baselines}.}\footnote{We note that finetuning separate classifier for semantic adherence and physical commonsense did not provide any additional benefits over a single classifier (\model{}) trained in a multi-task manner.} We assess auto-rater effectiveness by computing the ROC-AUC between human annotations and model judgments for videos generated from testing prompts.

% We will provide more details on the train and test set in \S \ref{sec:setup_dataset}.